\section{Related work}

\paragraph{Mobile and GUI agent benchmarks.} Interactive Android environments
with programmatic reward signals were established by
AndroidWorld~\citep{rawles2025androidworld}, which supplies tasks across real
applications on a live emulator rather than a static test set. MobileWorld
\citep{kong2026mobileworld}, the benchmark studied here, extends that setting
with agent--user interaction and MCP-augmented tasks and publishes a leaderboard
together with per-model trajectories. Our use of it is not adversarial: the
published trajectories are what made the zero-compute analysis of
Section~\ref{sec:zerocompute} possible at all, and few benchmarks release
enough for an outsider to do it.

\paragraph{Nondeterminism in language model inference.}
\citet{yuan2025nondeterminism} establish that low-precision inference is not
reproducible across serving configurations, tracing it to the non-associativity
of floating-point arithmetic and showing that batch size, GPU count and GPU
version alone can move accuracy by several points under greedy decoding. That
work explains the mechanism behind our Section~\ref{sec:determinism} and,
plausibly, behind Section~\ref{sec:endpoint}. The difference in setting matters:
their variables are controllable by whoever owns the hardware, whereas a
practitioner calling a hosted endpoint controls none of them and cannot observe
them either --- batch composition is a function of other customers' traffic.

\paragraph{Drift in hosted model services.} That a commercial model endpoint
changes behaviour over time is established. \citet{chen2024chatgpt} evaluated
the March and June 2023 releases of two hosted models across several task
families and found large swings --- identification of prime versus composite
numbers fell from 84\% to 51\% accuracy on one of them --- and argued for
continuous monitoring. The same group had earlier framed the general problem of
detecting shifts in machine learning APIs efficiently~\citep{chen2021apishifts}.

Our setting is narrower and, we think, harder to defend against. That work
compares \emph{labelled versions} of a service \emph{months apart}, and detects
the change by periodically re-running a whole evaluation suite. What we record
happened inside a single model identifier, at one provider, at one
quantization, within \emph{a single half-day}, with no version change announced or
announceable --- and it was caught not by re-running an evaluation but by a
cheap side signal, completion tokens per action, that costs one integer per
request to log. The first result tells a practitioner that a service will
change. The second tells them it can change \emph{during their experiment},
while every control they know about is still set, and that the standard check
will not show it.

\paragraph{Flaky tests.} Software engineering has a mature literature on tests
that pass and fail without a code change, including how practitioners
experience and manage them~\citep{gruber2022flakiness}. The analogy is close
and the vocabulary is useful, but the disanalogy is the point of this paper: a
flaky test is understood as a defect to be fixed, whereas run-to-run variance
in an environment-driven agent benchmark is a property of the measurement that
should be quantified and reported rather than eliminated. The software
engineering framing also assumes the system under test is the thing that
varies. Here the harness, the environment and the serving endpoint all vary,
and only one of them is the object of measurement.

\paragraph{Statistical reporting for evaluations.}
\citet{miller2024errorbars} sets out error-bar methodology for LLM evals,
\citet{bowyer2025clt} caution that CLT intervals understate uncertainty at
task counts like ours, and $\tau$-bench~\citep{yao2025taubench} reports
run-to-run inconsistency as a metric, pass\textasciicircum{}$k$. All three
hold the thing measured stationary while it is measured.
Section~\ref{sec:noise} is what their machinery yields here;
Section~\ref{sec:endpoint} is the case none of them models --- the instrument
moving mid-experiment.

\paragraph{Positioning.} These lines of work measure model or agent capability,
or the reliability of a test suite. This paper measures the instrument: how far
a benchmark's own output moves when its input is held as fixed as the available
controls allow, and what a reader would need recorded in order to notice when
it moves for reasons that have nothing to do with the agent.
